\section*{अध्याय \ref{ch:g12:equilibrium} --- रासायनिक साम्य: भागफल और स्थिरांक}

\begin{solution}{exo:g12:equilibrium:1}
(a) $Q = \dfrac{[\ce{NH3}]^2 (c^\circ)^2}{[\ce{N2}][\ce{H2}]^3}$;
(b) $Q = \dfrac{[\ce{CH3COO-}][\ce{H3O+}]}{[\ce{CH3COOH}]\,c^\circ}$ (पानी,
जो \omterm{def:g6:solutions-and-solubility:solute}{विलायक} है, छोड़ दिया जाता है);
(c) $Q = [\ce{Ca^{2+}}][\ce{CO3^{2-}}]/(c^\circ)^2$ (ठोस छोड़ दिया जाता है);
(d) $Q = \dfrac{[\ce{FeSCN^{2+}}]\,c^\circ}{[\ce{Fe^{3+}}][\ce{SCN-}]}$।
\end{solution}

\begin{solution}{exo:g12:equilibrium:2}
$\tau = 0.020/0.050 = 0.40$: पूर्ण नहीं।
\end{solution}

\begin{solution}{exo:g12:equilibrium:3}
$Q = 2 < K$: अग्र। $Q = 50 > K$: प्रतीप। $Q = K$: कोई परिवर्तन नहीं।
\end{solution}

\begin{solution}{exo:g12:equilibrium:4}
संरचना स्थिर है, पर अग्र और प्रतीप अभिक्रियाएँ समान दर पर
चलती रहती हैं: \omterm{def:g7:atoms-and-molecules:molecule}{अणु} दोनों दिशाओं में अभिक्रिया करते रहते हैं।
\end{solution}

\begin{solution}{exo:g12:equilibrium:5}
नहीं, $K$ अपरिवर्तित रहता है (\omterm{def:g10:reaction-progress-table:system}{अंतिम अवस्था} वही है); साम्य तक पहुँचने का
समय कम हो जाता है।
\end{solution}

\begin{solution}{exo:g12:equilibrium:6}
$\dfrac{x^2}{(2-x)^2} = 4$, इसलिए $\dfrac{x}{2-x} = 2$ और
$x = \qty{1.33}{mol}$: अम्ल और \omterm{def:g11:functional-groups:alcohol}{ऐल्कोहॉल} के \qty{0.67}{mol},
\omterm{def:g11:functional-groups:carboxylic-acid}{एस्टर} और पानी के \qty{1.33}{mol} (फिर से दो तिहाई)।
\end{solution}

\begin{solution}{exo:g12:equilibrium:7}
$Q = (0.5 \times 0.5)/(0.5 \times 0.5) = 1 < 4$: अग्र, और \omterm{def:g11:functional-groups:carboxylic-acid}{एस्टर}
बनता है।
\end{solution}

\begin{solution}{exo:g12:equilibrium:8}
अम्ल वाली ओर से लगभग \qty{0.52}{mol} और \omterm{def:g11:functional-groups:carboxylic-acid}{एस्टर} वाली ओर से
\qty{0.79}{mol}। लगभग \qty{100}{h} के बाद दोनों वक्र डैश वाली रेखा पर हैं:
साम्य।
\end{solution}

\begin{solution}{exo:g12:equilibrium:9}
$y$ mol जल-अपघटित होने पर: $\dfrac{y^2}{(1-y)^2} = \dfrac{1}{4}$, इसलिए
$\dfrac{y}{1-y} = \dfrac{1}{2}$, $y = \frac{1}{3}$: \omterm{def:g11:functional-groups:carboxylic-acid}{एस्टर} और पानी के
$\frac{2}{3}$ mol, अम्ल और \omterm{def:g11:functional-groups:alcohol}{ऐल्कोहॉल} के $\frac{1}{3}$ mol। एस्टरीकरण वाला ही
अंतिम \omterm{def:g6:pure-substances-and-mixtures:mixture}{मिश्रण}।
\end{solution}

\begin{solution}{exo:g12:equilibrium:10}
ठोस एक शुद्ध प्रावस्था है, जिसकी ``सांद्रता'' नहीं बदलती: उसे
$Q$ से बाहर रखा जाता है। और ठोस मिलाने से $Q$ नहीं बदलता, इसलिए वह
साम्य को नहीं खिसकाता (जब तक कुछ ठोस मौजूद है)।
\end{solution}

\begin{solution}{exo:g12:equilibrium:11}
$K$ गैस में कार्बन डाइऑक्साइड और पानी में कार्बन डाइऑक्साइड का अनुपात तय करता है।
खुलने पर पानी के ऊपर की गैस \omterm{def:g4:air-a-mixture-of-gases:air}{वायु} में निकल जाती है: गैस में कार्बन डाइऑक्साइड
घटती है, $Q < K$, और पानी से और गैस निकलती है, जब तक लगभग
कुछ न बचे।
\end{solution}

\begin{solution}{exo:g12:equilibrium:12}
साम्य पर: अम्ल और \omterm{def:g11:functional-groups:alcohol}{ऐल्कोहॉल} $\frac{1}{3}$, \omterm{def:g11:functional-groups:carboxylic-acid}{एस्टर} और पानी
$\frac{2}{3}$। आधा पानी हटाने पर $\frac{1}{3}$ बचता है:
$Q = \dfrac{(2/3)(1/3)}{(1/3)(1/3)} = 2 < 4$, अग्र। $y$ और
अभिक्रिया करने पर: $(2/3 + y)(1/3 + y) = 4(1/3 - y)^2$, अर्थात
$27y^2 - 33y + 2 = 0$, $y = 0.064$: \omterm{def:g11:functional-groups:carboxylic-acid}{एस्टर} बढ़कर \qty{0.73}{mol} हो जाता है।
\end{solution}

\begin{solution}{exo:g12:equilibrium:13}
उलटी अभिक्रिया के भागफल में अंश और हर
आपस में बदल जाते हैं: $Q' = 1/Q$, इसलिए साम्य पर $K' = 1/K$। जल-अपघटन:
$1/4 = 0.25$।
\end{solution}

\begin{solution}{exo:g12:equilibrium:14}
$Q = (2 \times 2)/(1 \times 1) = 4 = K$: \omterm{def:g6:pure-substances-and-mixtures:mixture}{मिश्रण} पहले से ही
साम्य में है; उसकी संरचना नहीं बदलती।
\end{solution}

\begin{solution}{exo:g12:equilibrium:15}
$n - 0.95 = 0.9025 / (4 \times 0.05) = 4.51$, इसलिए प्रति \omterm{def:g10:the-mole:amount}{मोल} अम्ल
$n = \qty{5.5}{mol}$ \omterm{def:g11:functional-groups:alcohol}{ऐल्कोहॉल}।
\end{solution}

\begin{solution}{pb:g12:equilibrium:1}
\textbf{1.} \ce{CH3-COOH + CH3-CH2-OH <=> CH3-COO-CH2-CH3 + H2O}।

\textbf{2.} वह पूर्ण नहीं है: वह एक साम्य पर रुक जाती है जहाँ दोनों
दिशाएँ चलती रहती हैं।

\textbf{3.} $Q = \dfrac{n_{\text{एस्टर}}\, n_{\text{पानी}}}{n_{\text{अम्ल}}\, n_{\text{ऐल्कोहॉल}}}$।

\textbf{4.} $x_{\max} = \qty{1}{mol}$; $x_f = \frac{2}{3}$ mol।

\textbf{5.} $\tau = 0.67$।

\textbf{6.} अम्ल और \omterm{def:g11:functional-groups:alcohol}{ऐल्कोहॉल} हर एक $\frac{1}{3}$ mol, \omterm{def:g11:functional-groups:carboxylic-acid}{एस्टर} और पानी हर एक
$\frac{2}{3}$ mol।

\textbf{7.} $K = \dfrac{(2/3)^2}{(1/3)^2} = 4$।

\textbf{8.} एक तिहाई जल-अपघटित: \omterm{def:g11:functional-groups:carboxylic-acid}{एस्टर} और पानी $\frac{2}{3}$, अम्ल
और \omterm{def:g11:functional-groups:alcohol}{ऐल्कोहॉल} $\frac{1}{3}$: वही \omterm{def:g6:pure-substances-and-mixtures:mixture}{मिश्रण}, $Q_{eq} = 4$।

\textbf{9.} यह उसी ताप पर वही अभिक्रिया है, और $K$
केवल ताप पर निर्भर करता है।

\textbf{10.} अम्ल $1 - x$; \omterm{def:g11:functional-groups:alcohol}{ऐल्कोहॉल} $3 - x$; \omterm{def:g11:functional-groups:carboxylic-acid}{एस्टर} $x$; पानी $x$।

\textbf{11.} $x^2 = 4(1 - x)(3 - x) = 4(3 - 4x + x^2)$, इसलिए
$3x^2 - 16x + 12 = 0$।

\textbf{12.} $x = (16 \pm \sqrt{256 - 144})/6 = (16 \pm 10.58)/6$: 0.903
या 4.43; केवल $x_f = \qty{0.903}{mol}$ 0 और 1 के बीच है।

\textbf{13.} \qty{90}{\percent}।

\textbf{14.} $M = \qty{88.0}{g/mol}$: $0.903 \times 88.0 = \qty{79.5}{g}$।

\textbf{15.} पानी का हर निष्कासन $Q$ का अंश घटाता है, जो
$K$ से नीचे गिर जाता है: निकाय फिर से आगे बढ़ता है।

\textbf{16.} अभिक्रिया तब तक चलती रहती जब तक अम्ल खप न जाए:
\qty{100}{\percent}।

\textbf{17.} एथेनॉल सस्ता है, और उसकी अधिकता \omterm{def:g11:functional-groups:carboxylic-acid}{एस्टर} से आसानी से अलग
(और पुनर्चक्रित) की जा सकती है।

\textbf{18.} केवल समय: जो अभिक्रिया लगभग कोई ऊष्मा नहीं देती, उसका
स्थिरांक ताप के साथ मुश्किल से बदलता है, और गरम करने से वह केवल
जल्दी साम्य तक पहुँचती है।

\textbf{19.} अम्ल का लगभग \qty{90}{\percent}।
\end{solution}
